\section{Results}
\subsection{Analytic sample}
The primary file contained 14,569 eligible player-day pairs from 50 players. Team A contributed 3,812 pairs in 2020 and 5,881 in 2021, while Team B contributed 2,393 and 2,483 pairs, respectively (Table~\ref{tab:sample}; Figure~\ref{fig:sample}). Across all pairs, both current-day and next-day readiness had a median of 7.0 (interquartile range 6.0--8.0). The median daily load was 350 session-RPE units (interquartile range 20--600), and 75.0\% of eligible pairs occurred on days with non-zero training load.

\begin{table}[htbp]
\centering
\caption{Analytic sample by team and year. Player counts are reported within team-year strata and are not additive across years.}
\label{tab:sample}
\scriptsize
\resizebox{\textwidth}{!}{%
\begin{tabular}{lccccc}
\toprule
Measure & Team A 2020 & Team A 2021 & Team B 2020 & Team B 2021 & Overall \\
\midrule
Players contributing eligible pairs & 23 & 27 & 14 & 23 & 50 \\
Eligible player-day pairs & 3,812 & 5,881 & 2,406 & 2,483 & 14,582 \\
Current readiness, median [IQR] & 7 [6, 8] & 7 [6, 8] & 7 [6, 8] & 7 [6, 8] & 7 [6, 8] \\
Next-day readiness, median [IQR] & 7 [6, 8] & 7 [6, 8] & 7 [6, 8] & 7 [6, 8] & 7 [6, 8] \\
Daily load (sRPE), median [IQR] & 360 [150, 630] & 320 [20, 560] & 360 [0, 560] & 360 [0, 630] & 350 [12, 600] \\
Pairs with non-zero daily load, n (\%) & 3,100 (81.3) & 4,415 (75.1) & 1,766 (73.4) & 1,657 (66.7) & 10,938 (75.0) \\
7-day readiness-reporting rate, median [IQR] & 1.00 [0.86, 1.00] & 1.00 [0.86, 1.00] & 1.00 [0.86, 1.00] & 0.86 [0.71, 1.00] & 1.00 [0.86, 1.00] \\
\bottomrule
\end{tabular}%
}
\end{table}


\begin{figure}[htbp]
\centering
\includegraphics[width=0.75\textwidth]{Figure2_analytic_sample.pdf}
\caption{Analytic sample by team and year. Bars show the number of eligible player-day pairs with complete current-day wellness data and observed next-day readiness. Labels report contributing players in each team-year subset.}
\label{fig:sample}
\end{figure}

\subsection{Primary incremental-value comparison}
Under the primary 28-observation sequential residual window, \Pone{} did not improve MAE over \Pzero{} in any evaluation setting (Table~\ref{tab:primary}; Figure~\ref{fig:incremental}). Within Team A, \Pzero{} and \Pone{} had MAEs of 0.722 and 0.738, respectively (\DeltaMAE{} +0.016, 95\% CI $-$0.001 to +0.036). Within Team B, MAE was 0.632 for \Pzero{} and 0.688 for \Pone{} (\DeltaMAE{} +0.056, +0.009 to +0.102). In cross-team evaluation, \DeltaMAE{} was +0.047 (+0.005 to +0.099) from Team A to Team B and +0.025 (+0.007 to +0.043) from Team B to Team A. Thus, three of four player-cluster bootstrap intervals excluded zero in the direction of higher error for \Pone{}, whereas the remaining within-Team-A comparison was close to null.

\begin{table}[htbp]
\centering
\caption{Primary temporal and two-team transportability results using a 28-observation sequential residual-history window. Positive \DeltaMAE{} values indicate higher mean absolute error for \Pone. Cross-team results use sequential player-specific residual adaptation based only on earlier observed target-team readiness reports.}
\label{tab:primary}
\scriptsize
\resizebox{\textwidth}{!}{%
\begin{tabular}{lrrrrrrr}
\toprule
Evaluation setting & Test pairs & Test players & \Pzero MAE & \Pone MAE & \Pone$-$\Pzero MAE (95\% CI) & \Pzero RMSE & \Pone RMSE \\
\midrule
Within Team A: 2020$\rightarrow$2021 & 5,881 & 27 & 0.722 & 0.738 & +0.016 ($-$0.001 to +0.036) & 1.044 & 1.046 \\
Within Team B: 2020$\rightarrow$2021 & 2,483 & 23 & 0.632 & 0.688 & +0.056 (+0.009 to +0.102) & 1.003 & 0.997 \\
Team A 2020$\rightarrow$Team B 2021 & 2,483 & 23 & 0.623 & 0.670 & +0.047 (+0.005 to +0.099) & 0.975 & 0.983 \\
Team B 2020$\rightarrow$Team A 2021 & 5,881 & 27 & 0.734 & 0.759 & +0.025 (+0.007 to +0.043) & 1.072 & 1.076 \\
\bottomrule
\end{tabular}%
}
\end{table}


\begin{figure}[htbp]
\centering
\includegraphics[width=0.90\textwidth]{Figure3_primary_incremental_value.pdf}
\caption{Primary incremental-value comparison. Points and horizontal intervals show player-cluster bootstrap estimates and 95\% confidence intervals for \DeltaMAE{} = MAE(\Pone) $-$ MAE(\Pzero) under the primary 28-observation sequential calibration window. Positive values indicate higher absolute error for \Pone.}
\label{fig:incremental}
\end{figure}

\subsection{Sensitivity analyses and calibration}
Changing the sequential residual window from 14 to 56 earlier eligible forecasts did not reverse the direction of the primary comparison (Figure~\ref{fig:window}; Supplementary Table~S1). Restricting evaluation to observations with 28 available player-specific residuals produced the same directional pattern. The reduced monitoring model \Ptwo{} also did not demonstrate incremental MAE improvement over \Pzero{} (Supplementary Table~S2).

Calibration slopes for both \Pzero{} and \Pone{} were below one in every scenario, indicating that predicted readiness values were somewhat too extreme relative to the observed variation under the stated linear calibration diagnostic (Figure~\ref{fig:calibration}; Supplementary Table~S3). Player-level \Pone{}--\Pzero{} MAE differences varied, but the median difference was positive in every scenario (Figure~\ref{fig:heterogeneity}).

\begin{figure}[htbp]
\centering
\includegraphics[width=0.90\textwidth]{Figure4_calibration_window_sensitivity.pdf}
\caption{Sensitivity to the residual-history window. Player-cluster bootstrap estimates and 95\% confidence intervals for 14-, 28-, and 56-observation windows. Positive values indicate higher absolute error for \Pone.}
\label{fig:window}
\end{figure}

\begin{figure}[htbp]
\centering
\includegraphics[width=0.90\textwidth]{Figure5_calibration_slope.pdf}
\caption{Calibration slope diagnostic. Calibration slopes were estimated from linear regressions of observed readiness on predicted readiness. The dashed reference is slope = 1.}
\label{fig:calibration}
\end{figure}

\begin{figure}[htbp]
\centering
\includegraphics[width=0.90\textwidth]{Figure6_player_level_incremental_error.pdf}
\caption{Player-level heterogeneity. Each point represents one player's \Pone{}--\Pzero{} MAE difference for the primary 28-observation analysis. Positive values indicate higher error for \Pone; black segments denote scenario-specific medians.}
\label{fig:heterogeneity}
\end{figure}
\FloatBarrier
